\section{Введение}

Обучение с подкреплением на данный момент является основным методом обучения агентов в условиях непрерывного принятия решений. Классические методы глубокого RL, такие как алгоритм DQN (Deep Q-Network), отлично это показывают. Однако и у него есть одна существенная проблема: инженеру необходимо заранее знать, какую архитектуру нейронной сети использовать для определенной задачи.

Обойти это ограничение можно с помощью метода нейроэволюции, в частности алгоритма NEAT (NeuroEvolution of Augmenting Topologies). Его основное отличие заключается в отказе от фиксированной архитектуры. Процесс обучения начинается с наименьшей топологии сети, которая усложняется только по мере необходимости. Это дает возможность оптимизировать структуру нейросети под нужды решения конкретной задачи.

В данной работе проводится сравнительный анализ классического глубокого Q-обучения и нейроэволюционного подхода.